% ==============================================================================
% core_step3_pcontrol.tex — SCI/SICE tutorial core deck, Step 3: P control and
% model-based design (exercises 5 and 6)
% SCI/SICE チュートリアル要点版 Step 3: P 制御とモデルに基づく設計（実習 5, 6）
%
% Gains are those of the sci2026 lesson code
% (firmware/workshop/lessons/lesson_05_p_control, lesson_06_modeling):
% exercise 5 = Kp 0.5 (roll/pitch), 2.0 (yaw); exercise 6 = zeta design.
% ここで使うゲインは sci2026 レッスンの実コードの値に合わせてある。
% ==============================================================================

% ------------------------------------------------------------------
\begin{frame}{これだけでは飛べない: フィードバックと内側ループ}
  \begin{center}
    \adjustbox{max width=0.92\textwidth, max height=0.50\textheight}{\documentclass[border=10pt]{standalone}
\usepackage{tikz}
\usetikzlibrary{arrows.meta, positioning, shapes.geometric, calc}

\begin{document}

\definecolor{blockcolor}{RGB}{0, 120, 200}
\definecolor{arrowcolor}{RGB}{80, 80, 80}
\definecolor{fbcolor}{RGB}{220, 50, 30}
\definecolor{sumcolor}{RGB}{40, 160, 40}
\begin{tikzpicture}[
    block/.style={
      rectangle, draw=blockcolor, line width=1.5pt,
      fill=blockcolor!8, rounded corners=3pt,
      minimum width=26mm, minimum height=14mm,
      font=\bfseries, text=blockcolor, align=center
    },
    % Summing-junction circle diameter is ~40% of block height (Block
    % Diagram Style convention), not 71% as before (10mm / 14mm block).
    % 加え合わせ点の円の直径はブロック高さの約4割（ブロック線図の流儀）。
    % 以前は 10mm / ブロック高さ14mm = 71% だった。
    sum/.style={
      circle, draw=sumcolor, line width=1.5pt,
      fill=sumcolor!8, minimum size=6mm,
      font=\normalsize\bfseries, text=sumcolor
    },
    arrow/.style={-{Stealth[length=3mm]}, line width=1.2pt, arrowcolor},
    lbl/.style={font=\small, above, text=arrowcolor},
  ]

  % Blocks
  \node[block] (target) at (0, 0) {Target\\Rate};
  \node[sum]   (sum)    at (3, 0) {$\Sigma$};
  \node[block] (kp)     at (6, 0) {P Control\\$K_p \times e$};
  \node[block] (mixer)  at (9.5, 0) {Motor\\Mixer};
  \node[block] (drone)  at (13, 0) {StampFly\\Drone};

  % Forward path
  \draw[arrow] (target) -- node[lbl] {$r$} (sum);
  \draw[arrow] (sum) -- node[lbl] {$e$} (kp);
  \draw[arrow] (kp) -- node[lbl] {$u$} (mixer);
  \draw[arrow] (mixer) -- (drone);

  % Output line with feedback tap point
  \draw[arrow] (drone.east) -- ++(2.5, 0) node[lbl, right] {response};
  \coordinate (fb_tap) at ($(drone.east)+(1.2, 0)$);
  \fill[arrowcolor] (fb_tap) circle (2.5pt);

  % Feedback path (branches from output line, not from inside drone)
  \node[block, fill=fbcolor!8, draw=fbcolor, text=fbcolor]
    (gyro) at (8, -3) {Gyro\\Sensor};
  \draw[arrow, fbcolor] (fb_tap) -- (fb_tap |- gyro) -- (gyro);
  \draw[arrow, fbcolor] (gyro) -- (3, -3) -- (sum);

  % Sum signs: "+" above the left-entering arrow, "-" to the RIGHT of the
  % bottom-entering feedback arrow (Block Diagram Style convention).
  % 加え合わせ点の符号: 「+」は左から入る矢印の上、「-」は下から入る
  % フィードバック矢印の右側に置く（ブロック線図の流儀）。
  % Fix: the "-" was at 210 deg (left of the upward arrow) with anchor
  % "north east" (glyph extending further left) -- it sat on the WRONG
  % (left) side of the arrow. Moved to 330 deg (right of the arrow) with
  % anchor "north west" (glyph extending right of the point).
  % 修正: 以前「-」は210度（上向き矢印の左側）・anchor=north east
  % （グリフがさらに左に伸びる）で、矢印の誤った（左）側にあった。
  % 330度（矢印の右側）・anchor=north west（グリフが右に伸びる）へ移動。
  \node[font=\scriptsize, sumcolor, anchor=south east] at ($(sum.north)+(150:0.3)$) {$+$};
  \node[font=\scriptsize, fbcolor, anchor=north west] at ($(sum.south)+(330:0.3)$) {$-$};

  % Labels
  \node[font=\small, arrowcolor, above] at ($(drone.east)+(1.2, 0.5)$) {angular rate};
  \node[font=\small, fbcolor, below] at (8, -3.8) {feedback};

\end{tikzpicture}
\end{document}
}
  \end{center}

  \vspace{-0.2em}
  \begin{itemize}\footnotesize\setlength{\itemsep}{1pt}
    \item 実習 4 の開ループ（スティック $\to$ ミキサ $\to$ モータ）では、モータの個体差、機体の非対称（推力軸・重心のずれ）、電池電圧の低下に何の補正もかからない
    \item そこで角速度をジャイロで測り、目標との差（誤差 $e=r-\omega$）に比例ゲイン $K_p$ を掛けてミキサへ戻す。これを\textbf{内側ループ（レート制御）}と呼ぶ
    \item 姿勢角の制御はこの外側に重なる（STABILIZE 以上のモード）。今日は内側ループだけを扱う
  \end{itemize}
\end{frame}

% ------------------------------------------------------------------
\begin{frame}[fragile]{実習 5: レート P 制御で飛ばす}
  \exgoal{ジャイロでフィードバックをかけ、制御の効果を見る（$K_p$ は勘で決めた値）}
  \begin{columns}[T]
    \begin{column}{0.44\textwidth}
      {\footnotesize\textbf{\textcolor{sfblue}{手順}}}\par\vspace{-2pt}
      \footnotesize
      \begin{enumerate}\setlength{\itemsep}{0pt}
        \item \code{sf lesson switch sci2026:5}
        \item \code{sf lesson edit} で TODO を埋めて保存
        \item \code{sf lesson build}（ここまで各自）
        \item 飛行は講師が完成コードで行う
      \end{enumerate}
    \end{column}
    \begin{column}{0.54\textwidth}
      \begin{sflisting}[basicstyle=\ttfamily\scriptsize, moredelim={[is][\bfseries\color{sfred}]{<<}{>>}}]{user\_code.cpp（赤字を書く）}
Kp_rp = 0.5f; Kp_yaw = 2.0f;
<<float rt = ws::rc_roll() * rate_max_rp;>>
<<float re = rt - ws::gyro_x();>>
<<float ro = Kp_rp * re;>>
// pitch, yaw: same form
<<ws::set_rate_target(rt, pt, yt);>>
ws::motor_mixer(ws::rc_throttle(),
                ro, po, yo);
      \end{sflisting}
    \end{column}
  \end{columns}
  \exresult{スティックに機体が追従し、わずかな振動が見える}
  \exhelptodo{5}
\end{frame}

% ------------------------------------------------------------------
\begin{frame}{プラントの伝達関数}
  \begin{center}
    \vspace{-2mm}
    \adjustbox{max width=0.80\textwidth, max height=0.34\textheight}{\documentclass[border=10pt]{standalone}
\usepackage{tikz}
\usepackage{amsmath}
\usetikzlibrary{arrows.meta, positioning, shapes.geometric, calc}

\begin{document}

\definecolor{blockcolor}{RGB}{0, 120, 200}
\definecolor{arrowcolor}{RGB}{80, 80, 80}
\definecolor{plantcolor}{RGB}{40, 160, 40}
\definecolor{sumcolor}{RGB}{120, 80, 150}
\begin{tikzpicture}[
    block/.style={
      rectangle, draw=#1, line width=1.5pt,
      fill=#1!8, rounded corners=3pt,
      minimum width=38mm, minimum height=16mm,
      font=\bfseries, text=#1, align=center
    },
    arrow/.style={-{Stealth[length=3mm]}, line width=1.2pt, arrowcolor},
    lbl/.style={font=\small, above, text=arrowcolor},
  ]

  % Input
  \node[font=\bfseries, arrowcolor] (input) at (-1.5, 0) {$u$};

  % Motor block (actuator + 1st-order lag)
  \node[block=blockcolor] (motor) at (3.5, 0)
    {Mixer + Motor\\[2pt]$\dfrac{K_m}{\tau_m s + 1}$};

  % Body block (rigid-body integrator)
  \node[block=plantcolor] (body) at (10, 0)
    {Body\\[2pt]$\dfrac{1}{I \cdot s}$};

  % Output
  \node[font=\bfseries, arrowcolor] (output) at (14.5, 0) {$\omega$};

  % Arrows
  \draw[arrow] (input) -- (motor);
  \draw[arrow] (motor) -- node[lbl] {$\tau$} (body);
  \draw[arrow] (body)  -- (output);

  % Labels
  \node[font=\scriptsize, blockcolor, below=2mm] at (motor.south)
    {actuator dynamics};
  \node[font=\scriptsize, plantcolor, below=2mm] at (body.south)
    {rigid-body rotation};
  \node[font=\scriptsize, arrowcolor, above=1mm] at (input.north)
    {control input};
  \node[font=\scriptsize, arrowcolor, above=1mm] at (output.north)
    {angular rate};

  % Collapsed transfer function
  \node[font=\normalsize, arrowcolor, anchor=north] at (6.75, -2.8) {%
    $G_p(s) = \dfrac{K}{s(\tau_m s + 1)},\qquad K = \dfrac{K_m}{I}$};

\end{tikzpicture}
\end{document}
}
  \end{center}

  \begin{exampleblock}{ホバー近傍で線形化すると 2 ブロックの直列になる}\footnotesize
    \begin{itemize}\setlength{\itemsep}{1pt}
      \item ミキサ + モータ: Duty $\to$ トルク、1 次遅れ $K_m/(\tau_m s+1)$
      \item 機体: トルク $\to$ 角速度、積分器 $1/(Is)$（$I$: 慣性モーメント）
      \item 直列にすると $G_p(s)=\dfrac{K}{s(\tau_m s+1)}$、 $K=K_m/I$
    \end{itemize}
  \end{exampleblock}
\end{frame}

% ------------------------------------------------------------------
\begin{frame}{実測パラメータ}
  \begin{table}
    \centering\footnotesize
    \begin{tabular}{llccc}
      \hline
      \textbf{パラメータ} & \textbf{記号} & \textbf{Roll} & \textbf{Pitch} & \textbf{Yaw} \\
      \hline
      慣性モーメント & $I$ & 9.16e-6 & 13.3e-6 & 20.4e-6\,kg$\cdot$m$^2$ \\
      トルクゲイン   & $K_m$ & 9.3e-4 & 9.3e-4 & 1.67e-4\,N$\cdot$m \\
      \hline
      \textbf{プラントゲイン} & $K=K_m/I$ & \textbf{102} & \textbf{70} & \textbf{$\approx$8}\,rad/s$^2$ \\
      \hline
    \end{tabular}
  \end{table}

  \vspace{0.4em}

  \begin{itemize}\footnotesize\setlength{\itemsep}{2pt}
    \item モータ時定数 $\tau_m$ は 0.02\,s を使う（3 軸共通。実測諸元からの計算値は 0.0163\,s）
    \item Yaw の $K$ は計算で 8.2、実習コードは 8.0 を使う。Roll/Pitch より 1 桁小さい（ヨートルクを出しにくい）ので、次ページの設計式では軸ごとに必要な $K_p$ が大きく違う
  \end{itemize}
\end{frame}

% ------------------------------------------------------------------
\begin{frame}{閉ループの 2 次系と、$\zeta$ から $K_p$ を決める式}
  \begin{columns}[c]
    \begin{column}{0.40\textwidth}
      \centering
      \adjustbox{max width=\columnwidth, max height=0.78\textheight}{\documentclass[border=10pt]{standalone}
\usepackage{tikz}
\usetikzlibrary{arrows.meta}

\begin{document}

\definecolor{refcolor}{RGB}{150, 150, 150}
\definecolor{labelcolor}{RGB}{80, 80, 80}
\definecolor{gridcolor}{RGB}{220, 220, 220}
% Three zeta colors
\definecolor{z05color}{RGB}{220, 50, 30}   % Red - underdamped
\definecolor{z07color}{RGB}{0, 120, 200}   % Blue - optimal
\definecolor{z10color}{RGB}{40, 160, 40}   % Green - critically damped
\begin{tikzpicture}[
    axis/.style={-{Stealth[length=2.5mm]}, line width=0.8pt, labelcolor},
    ref/.style={refcolor, line width=1pt, dashed},
    curve/.style={line width=2pt},
  ]

  % Grid
  \foreach \x in {0, 1, ..., 8} {
    \draw[gridcolor] (\x, 0) -- (\x, 5);
  }
  \foreach \y in {0, 1, ..., 5} {
    \draw[gridcolor] (0, \y) -- (8, \y);
  }

  % Axes
  \draw[axis] (-0.3, 0) -- (8.5, 0)
    node[right, font=\small\bfseries] {$t$};
  \draw[axis] (0, -0.3) -- (0, 5.5)
    node[above, font=\small\bfseries] {$\omega$};

  % Reference line at y=3.5 (steady state)
  \pgfmathsetmacro{\A}{3.5}
  \draw[ref] (0, \A) -- (8, \A);
  \node[font=\scriptsize, refcolor, left] at (-0.1, \A) {target};

  % ζ = 0.5 (underdamped, ωn=10)
  % y(t) = A*[1 - exp(-ζωn t)/√(1-ζ²) * sin(ωd t + φ)]
  % ζ=0.5, ωn=10, ωd=8.66, φ=60°
  \pgfmathsetmacro{\zetaA}{0.5}
  \pgfmathsetmacro{\wnA}{10}
  \pgfmathsetmacro{\zwA}{\zetaA*\wnA}         % 5.0
  \pgfmathsetmacro{\wdAdeg}{\wnA*sqrt(1-\zetaA*\zetaA)*180/3.14159} % ≈496°/s
  \pgfmathsetmacro{\phiAdeg}{acos(\zetaA)}     % 60°
  \pgfmathsetmacro{\invA}{1/sqrt(1-\zetaA*\zetaA)} % 1.155
  \draw[curve, z05color] plot[domain=0:2.0, samples=200]
    ({\x*4}, {\A*(1 - \invA*exp(-\zwA*\x)*sin(\wdAdeg*\x + \phiAdeg))});

  % ζ = 0.7 (slightly underdamped, ωn=10)
  \pgfmathsetmacro{\zetaB}{0.7}
  \pgfmathsetmacro{\wnB}{10}
  \pgfmathsetmacro{\zwB}{\zetaB*\wnB}         % 7.0
  \pgfmathsetmacro{\wdBdeg}{\wnB*sqrt(1-\zetaB*\zetaB)*180/3.14159} % ≈409°/s
  \pgfmathsetmacro{\phiBdeg}{acos(\zetaB)}     % 45.6°
  \pgfmathsetmacro{\invB}{1/sqrt(1-\zetaB*\zetaB)} % 1.400
  \draw[curve, z07color] plot[domain=0:2.0, samples=200]
    ({\x*4}, {\A*(1 - \invB*exp(-\zwB*\x)*sin(\wdBdeg*\x + \phiBdeg))});

  % ζ = 1.0 (critically damped, ωn=10)
  % y(t) = A*[1 - (1 + ωn t) exp(-ωn t)]
  \pgfmathsetmacro{\wnC}{10}
  \draw[curve, z10color] plot[domain=0:2.0, samples=200]
    ({\x*4}, {\A*(1 - (1 + \wnC*\x)*exp(-\wnC*\x))});

  % Legend
  \node[font=\small\bfseries, z05color, right] at (5.2, 4.8)
    {$\zeta = 0.5$};
  \node[font=\small\bfseries, z07color, right] at (5.2, 4.2)
    {$\zeta = 0.7$};
  \node[font=\small\bfseries, z10color, right] at (5.2, 3.0)
    {$\zeta = 1.0$};

  % Overshoot annotation for ζ=0.5
  \pgfmathsetmacro{\peakY}{\A*(1+exp(-3.14159*\zetaA/sqrt(1-\zetaA*\zetaA)))}
  \draw[z05color, line width=0.6pt, dashed]
    (1.45, \A) -- (1.45, \peakY);
  \node[font=\tiny, z05color, above] at (1.45, \peakY) {16\%};

\end{tikzpicture}
\end{document}
}
    \end{column}
    \begin{column}{0.58\textwidth}
      \footnotesize
      \setlength{\abovedisplayskip}{3pt}\setlength{\belowdisplayskip}{3pt}%
      P 制御の閉ループは 2 次系になる（$\omega_n$: 固有角周波数 [rad/s]、$\zeta$: 減衰比）:
      \[
        G_{cl}(s)=\frac{K_pK}{\tau_m s^2+s+K_pK}
      \]
      \[
        \omega_n=\sqrt{\frac{K_pK}{\tau_m}},\quad
        \zeta=\frac{1}{2\sqrt{K_pK\tau_m}}
      \]
      狙う $\zeta$ から $K_p$ を逆算する設計式:
      \[
        \boxed{K_p=\frac{1}{4\zeta^2K\tau_m}}
      \]
      オーバーシュート量 $\approx e^{-\pi\zeta/\sqrt{1-\zeta^2}}$（$\zeta=0.5\to16\%$、$\zeta=0.7\to5\%$、$\zeta\ge1\to0\%$）
    \end{column}
  \end{columns}
\end{frame}

% ------------------------------------------------------------------
\begin{frame}[fragile]{実習 6: 設計式で $K_p$ を計算して飛ばす}
  \exgoal{軸ごとに設計式から $K_p$ を求め、$\zeta$ を変えたときの違いを見る}
  \begin{columns}[T]
    \begin{column}{0.47\textwidth}
      {\footnotesize\textbf{\textcolor{sfblue}{手順}}}\par\vspace{-2pt}
      \footnotesize
      \begin{enumerate}\setlength{\itemsep}{0pt}
        \item \code{sf lesson switch sci2026:6}
        \item \code{sf lesson edit} で \code{Kp\_*} を設計式で計算して保存
        \item \code{sf lesson build}（ここまで各自）
        \item 飛行は講師が $\zeta=0.7\to0.5\to1.0$ の順に変えて行う
      \end{enumerate}
    \end{column}
    \begin{column}{0.51\textwidth}
      \begin{sflisting}[basicstyle=\ttfamily\scriptsize, moredelim={[is][\bfseries\color{sfred}]{<<}{>>}}]{user\_code.cpp（赤字を書く）}
tau_m = 0.02f;  K_roll = 102.0f;
K_pitch = 70.0f;  K_yaw = 8.0f;
float zeta = 0.7f;
<<float c = 1.0f / (4*zeta*zeta * tau_m);>>
<<float Kp_roll  = c / K_roll;>>
<<float Kp_pitch = c / K_pitch;>>
<<float Kp_yaw   = c / K_yaw;>>
      \end{sflisting}
    \end{column}
  \end{columns}
  \exresult{$\zeta=0.5$ は行き過ぎが大きく、$\zeta=1.0$ は反応が鈍い}
  \exhelptodo{6}
\end{frame}

% ------------------------------------------------------------------
\begin{frame}{実習 6 の結果: 軸ごとに $K_p$ は違う}
  \begin{table}
    \centering\footnotesize
    \begin{tabular}{lcccccc}
      \hline
      & $K$ & \textbf{実習 5 の $K_p$} & \textbf{その $\zeta$} & $\zeta=0.5$ の $K_p$ & $\zeta=0.7$ の $K_p$ & $\zeta=1.0$ の $K_p$ \\
      \hline
      Roll  & 102 & 0.5 & 0.50 & 0.49 & 0.25 & 0.12 \\
      Pitch &  70 & 0.5 & 0.60 & 0.71 & 0.36 & 0.18 \\
      Yaw   & 8.0 & 2.0 & 0.88 & 6.25 & 3.19 & 1.56 \\
      \hline
    \end{tabular}
  \end{table}

  \vspace{0.3em}

  \begin{exampleblock}{読み取り}\footnotesize
    実習 5 の勘の値は、Roll で $\zeta\approx0.5$、Pitch で $0.6$（やや振動的）、Yaw は $K_p=2.0$ で $\zeta\approx0.88$ だった。$K$ が軸ごとに 1 桁違うため、同じ $K_p$ を 3 軸に使い回せない。モデルがあれば、勘の値の良し悪しを数値で言える
  \end{exampleblock}
\end{frame}
